\documentclass[11pt,a4paper]{article}
\usepackage[utf8]{inputenc}
\usepackage[spanish]{babel}
\usepackage{amsmath}
\usepackage{amsfonts}
\usepackage{amssymb}
\usepackage{listings}
\usepackage{xcolor}
\usepackage{graphicx}
\usepackage{geometry}
\usepackage{hyperref}
\usepackage{tcolorbox}

\geometry{margin=2.5cm}

% Configuracion de colores para codigo
\definecolor{codegreen}{rgb}{0,0.6,0}
\definecolor{codegray}{rgb}{0.5,0.5,0.5}
\definecolor{codepurple}{rgb}{0.58,0,0.82}
\definecolor{backcolour}{rgb}{0.95,0.95,0.92}
\definecolor{codeblue}{rgb}{0.1,0.1,0.8}

\lstset{
    backgroundcolor=\color{backcolour},   
    commentstyle=\color{codegreen},
    keywordstyle=\color{magenta}\bfseries,
    numberstyle=\tiny\color{codegray},
    stringstyle=\color{codepurple},
    basicstyle=\ttfamily\footnotesize,
    breakatwhitespace=false,         
    breaklines=true,                 
    captionpos=b,                    
    keepspaces=true,                 
    numbers=left,                    
    numbersep=5pt,                  
    showspaces=false,                
    showstringspaces=false,
    showtabs=false,                  
    tabsize=2,
    language=Python,
    morekeywords={torch, nn, F, np, librosa, argparse}
}

\title{\textbf{Manual de Ingeniería: Refactorización IDEAW}\\ \large Guía Maestra de Implementación y Teoría}
\author{Documentación Técnica para Replicación}
\date{\today}

\begin{document}

\maketitle

\begin{tcolorbox}[colback=blue!5!white,colframe=blue!75!black,title=Objetivo del Proyecto]
Modificar el sistema de marcado de agua IDEAW para operar exclusivamente en la magnitud de las frecuencias bajas (0-100 bins), garantizando la transparencia auditiva total al preservar la fase original del audio.
\end{tcolorbox}

\tableofcontents
\newpage

\section{Teoría de Señales y Psicoacústica}

\subsection{¿Por qué Bajas Frecuencias?}
El oído humano es menos sensible a pequeños cambios de amplitud en frecuencias muy bajas o muy altas. IDEAW aprovecha esto. Al limitar el marcado a los primeros 100 bins (bajos), nos aseguramos de que la energía de la marca de agua se mezcle con los ritmos y bajos potentes, haciéndola casi indetectable al oído.

\subsection{El Poder de la Fase}
La fase contiene la información de "cuándo" ocurre cada frecuencia. Si la fase se altera, aunque sea un poco, el audio puede sonar "borroso" o "metálico". 
Al usar la fórmula:
$$Z_{wm} = |A_{wm}| \cdot e^{j\phi_{orig}}$$
Estamos creando un nuevo audio que tiene el volumen modificado ($|A_{wm}|$) pero mantiene la sincronización exacta del audio original ($\phi_{orig}$).

\section{Explicación Detallada del Código (\texttt{models/ideaw.py})}

\subsection{Bloque 1: Separación de Componentes}
\begin{lstlisting}[language=Python]
audio_stft = self.stft(audio) 
audio_complex = torch.view_as_complex(audio_stft)
audio_mag = torch.abs(audio_complex)
audio_phase = torch.angle(audio_complex)
\end{lstlisting}
\textbf{Explicación:} 
\begin{itemize}
    \item \texttt{stft}: Convierte la onda en una matriz de números.
    \item \texttt{view\_as\_complex}: PyTorch guarda la parte Real e Imag como dos números separados. Esto los junta en un solo número complejo matemático.
    \item \texttt{abs} y \texttt{angle}: Son funciones estándar para obtener el tamaño (volumen) y el ángulo (fase) de esos números complejos.
\end{itemize}

\subsection{Bloque 2: El truco de las dimensiones (\texttt{unsqueeze})}
\begin{lstlisting}[language=Python]
audio_bajos_in = audio_bajos_mag.unsqueeze(-1)
\end{lstlisting}
\textbf{Explicación:} 
Las redes neuronales convolucionales (CNN) esperan datos en 4 dimensiones: \texttt{[Batch, Canal, Altura, Ancho]}. Como al extraer la magnitud nos quedamos con una matriz plana, \texttt{unsqueeze(-1)} añade una dimensión "falsa" al final para que la red crea que hay 1 canal de color (como una imagen en blanco y negro).

\subsection{Bloque 3: Reconstrucción Cartesiana}
\begin{lstlisting}[language=Python]
wm_audio_complex = wm_mag_completa * torch.exp(1j * audio_phase)
wm_audio_stft = torch.view_as_real(wm_audio_complex)
\end{lstlisting}
\textbf{Explicación:} 
\begin{itemize}
    \item \texttt{torch.exp(1j * phase)}: Es la implementación en código de la fórmula de Euler. Crea un número complejo de magnitud 1 con el ángulo de la fase original.
    \item Al multiplicarlo por nuestra magnitud marcada, "estiramos" ese número hasta el volumen que queremos.
    \item \texttt{view\_as\_real}: Desarma el número complejo otra vez en partes Real e Imaginaria porque la función \texttt{istft} no sabe leer números complejos directamente.
\end{itemize}

\section{Flujo de Trabajo en \texttt{embed\_extract.py}}

\subsection{Sistema de Argumentos (\texttt{argparse})}
Añadimos flexibilidad para que no tengas que editar el código cada vez:
\begin{itemize}
    \item \texttt{--input}: Puede recibir un archivo \texttt{.wav} o una carpeta.
    \item \texttt{--output\_dir}: Donde se guardarán los resultados.
    \item \texttt{--checkpoint}: Permite cargar los pesos entrenados del modelo.
\end{itemize}

\subsection{La función \texttt{process\_file}}
Esta función es el "cerebro" de la ejecución:
\begin{enumerate}
    \item \textbf{Generación de marca:} Crea un mensaje binario al azar (ceros y unos).
    \item \textbf{Embedding:} Llama al modelo para esconder la marca.
    \item \textbf{Guardado:} Escribe el nuevo archivo de audio.
    \item \textbf{Extracción:} Inmediatamente intenta leer la marca del audio que acaba de crear para decirte si funcionó (ACC).
\end{enumerate}

\section{Cómo hacer que funcione (Paso a Paso)}

\begin{tcolorbox}[colback=green!5!white,colframe=green!75!black,title=Guía de Configuración]
\begin{enumerate}
    \item \textbf{Entorno:} Asegúrate de estar en el directorio raíz del proyecto.
    \item \textbf{Activación:} Ejecuta \texttt{.venv\textbackslash Scripts\textbackslash activate}.
    \item \textbf{Verificación:} Ejecuta \texttt{python --version}. Debe ser Python 3.x.
    \item \textbf{Ejecución Simple:} \texttt{python embed\_extract.py} (esto usará el audio por defecto).
    \item \textbf{Ejecución Masiva:} Pon tus audios en una carpeta y usa el comando \texttt{--input}.
\end{enumerate}
\end{tcolorbox}

\section{Preguntas Frecuentes (FAQ)}
\textbf{¿Por qué el ACC es bajo (0.5 - 0.6)?}\\
Porque el modelo está usando pesos al azar. Para que el ACC suba a 1.0, debes entrenar el modelo usando \texttt{train.py} con esta nueva estructura de 1 canal.

\textbf{¿Qué pasa si mi audio dura más de 5 segundos?}\\
El script lo divide en trozos de 16,000 muestras (1 segundo) y procesa cada trozo por separado automáticamente.

\end{document}
